\section{格点系综与色散关系}
\label{sec:simulate}
我们在 $N_f=2$ 非微扰改进 Wilson 费米子加 Wilson 规范作用、由 QCDSF
和 RQCD 生成的 200 个有效去关联构型上研究新的动量涂抹方法。这构成
文献~\cite{Bali:2014gha} 系综 IV 的一个子集。注意在
强子结构研究中我们通常在约 2000 个构型上采用多个源~\cite{Bali:2014gha}，即
目前的统计量是很有限的。尽管如此，
我们将报告高达 $2.8\units{GeV}$ 动量的有意义信号。这个 $32^3\times 64$ 系综的参数值为
$\beta=5.29$、$\kappa=0.13632$，对应逆格距
$a^{-1}\approx 2.76\units{GeV}$ 和（有限体积）
$\pi$ 介子质量 $m_{\pi}\approx 295\units{MeV}$。这给出空间尺度
$L\approx 2.29\units{fm}\approx 3.42/m_{\pi}$。
注意在这个格距下实现物理 $\pi$ 介子质量同时保持 $L\gtrsim 3.5/m_{\pi}$ 要求 $L>70a$。由于我们的目标本来就是超过强子质量的动量，我们不
预期向更小 $\pi$ 介子质量过渡会引起结果的定性变化，
故选择了当前参数作为足够现实但仍可负担的折中。

在每个空间方向点数为偶数 $N=L/a$ 的有限立方格点上，动量只能取离散值
\begin{equation}
\label{quantize}
\vect{p}=\frac{2\pi}{L}\vect{P}\,,\quad
P_i\in\left\{-\frac{L}{2a}+1,-\frac{L}{2a}+2,\cdots, \frac{L}{2a}\right\}\,.
\end{equation}
这意味着最小的非平凡 $|\vect{P}|=|(1,0,0)|$ 给出
动量 $|\vect{p}|\approx 0.54\units{GeV}$，而
$\vect{P}=(3,3,3)$ 对应于 $|\vect{p}|\approx 2.82\units{GeV}$。

$\pi$ 介子与核子质量已在文献~\cite{Bali:2014nma} 中以高统计量确定，为
\begin{equation}
\label{eq:masses}
am_{\pi}=0.10675(52)\,,\quad am_N=0.3855(45)\,.
\end{equation}
我们将利用这些参考值，把我们的 $\pi$ 介子与核子基态能量与连续及格点
色散关系的预期进行比较。
连续色散关系读作
\begin{equation}
\label{eq:disperse}
aE=a\sqrt{m^2+\vect{p}^2}\,.
\end{equation}
此外，我们把 $\pi$ 介子能量与自由朴素离散化标量粒子的
格点色散关系
\begin{equation}
\label{eq:dispersepi}
\cosh(aE_{\pi})=\cosh(am_{\pi})+\frac{a^2}{2}\hat{\vect{p}}^2\,,\quad
\end{equation}
比较，把核子能量与 Wilson 参数 $r=1$ 的自由 Wilson 费米子的预期依赖（参见例如
文献~\cite{ElKhadra:1996mp}）
\begin{equation}
\label{eq:disperseN}
\cosh(aE_N)=1+\frac{\left(e^{am_N}-1+a^2\hat{\vect{p}}^2/2\right)^2+a^2\overline{\vect{p}}^2}{2\left(e^{am_N}+a^2\hat{\vect{p}}^2/2\right)}\,.
\end{equation}
比较。以上我们使用了标准缩写
\begin{equation}
\hat{p}_{\mu}=\frac{2}{a}\sin\left(\frac{ap_{\mu}}{2}\right)\,,\quad
\overline{p}_{\mu}=\frac{1}{a}\sin(ap_{\mu})\,,
\end{equation}
及 $\hat{\vect{p}}^2=\sum_j\hat{p}_j^2$、
$\overline{\vect{p}}^2=\sum_j\overline{p}_j^2$。
我们指出朴素传播子中的质量参数 $m_0$
（例如标量情形：$1/(m_0^2+\hat{p}_{\mu}\hat{p}_{\mu})$）
与静止系能量相差格点人为效应。
上面的质量定义为满足 $m=E(\vect{0})$，理应如此。
它们转换到
$m_0$ 的公式分别给定为
$m_{\pi}=2a^{-1}\sinh (am_0/2)$ 与 $m_N=a^{-1}\ln(1+am_0)$。

\begin{figure}[ht]
\centerline{\includegraphics[width=0.95\textwidth]{smear3d.pdf}}
\caption{涂抹密度轮廓 $\rho(\vect{x})$
（式~\protect\eqref{density}）在 $x_1$-$x_2$ 平面内的截面。
颜色编码相位式~\protect\eqref{phase}，
动量涂抹参数 $\vect{k}$ 的方向为 $(1,1,0)$。
左上：自由场情形。右上：APE 涂抹规范链接。
左下：原始规范链接。右下：带附加 boost 因子 $\gamma=5.3$ 的 APE 涂抹链接。}
\label{fig:smearing}
\end{figure}

显然，连续色散关系在大动量处应当被破坏。在这种情形下，我们也并不预期适用于自由点状粒子的
格点色散关系能准确描述数据。不过，
连续公式与格点公式之差可以用来朴素地估计预期的格点人为效应的大小。
